% Root's ultralimit construction, independently checked and written by the Part F agent.
\subsection{A discontinuous action with no cyclotomic fixed pro-point}

\begin{claim}{D.12}{Proved}
For $U=\mathbb P^1_\Q\setminus\{0,1,\infty\}$, put
$E=\pi_1(U)$ and $\Pi=\pi_1(U_{\overline\Q})$, and fix a
decomposition embedding $G_{\Q_3}\hookrightarrow G_\Q$.
There is an abstract, provably discontinuous section
\[
                         s:G_\Q\longrightarrow E
\]
whose restriction to $G_{\Q_3}$ is geometrically conjugate
to the continuous point section at the explicit transcendental
point
\[
                    a=1+\sum_{r=1}^{\infty}3^{r!}
                       \in U(\Q_3).
\]
It induces an abstract semilinear action on the universal
pro-cover and its completed proper levels, normalizing the
original geometric deck group $\Pi$.
\end{claim}
\begin{proof}
First verify transcendence, rather than assume a generic local
point. Put $A_N=1+\sum_{r=1}^N3^{r!}$. Then
\[
 v_3(a-A_N)=(N+1)!,\qquad
                   1<A_N<2\cdot3^{N!}.
\]
If a nonzero $P\in\Z[T]$ of degree $d$ vanished at $a$,
then $3^{(N+1)!}$ would divide $P(A_N)$. The latter integer
is nonzero for all sufficiently large $N$, since the $A_N$
are distinct. On the other hand
$|P(A_N)|_\infty\leq C(2\cdot3^{N!})^d$ for a fixed
$C$. This contradicts $(N+1)!>dN!+O(1)$. Thus $a$ is
transcendental over $\Q$. It is a $3$-adic unit, and
$v_3(a-1)=1$, so it is an interior point of $U$.

Construct positive integers $x_n\geq2$ recursively. Choose a
prime $q_n\ne3$ which divides none of $x_1,\ldots,x_{n-1}$
and impose
\[
 \begin{split}
 x_n&\equiv a\pmod{3^n},\\
 x_n&\equiv q_n\pmod{q_n^2},\\
 x_n&\equiv1\pmod{q_i}\quad(1\leq i<n).
 \end{split}
\]
The primes $q_n$ are necessarily distinct, and the Chinese
remainder theorem gives such an integer. One may always choose
the least permitted prime and the least positive CRT representative.
That representative is neither zero nor one because of its
congruence modulo $q_n^2$. The first congruence is effective:
only terms $3^{r!}$ with $r!<n$ contribute. We have
\[
           x_n\longrightarrow a\text{ in }\Q_3,
          \qquad v_{q_i}(x_j)=\delta_{ij}\quad(i,j\geq1).
\]
For $j<i$ the valuation assertion is the choice of $q_i$;
for $j>i$ it is the third congruence. In particular the
square classes of the $x_n$ are independent. The exact
twelve-term CRT and valuation certificate is in
\texttt{code/F\_D5\_ultralimit\_output.json}; its role is
only to check this finite arithmetic input, not the
ultrafilter or the fixed-point argument.

Choose any continuous representatives $s_n$ of the global
point sections at $x_n$. Fix a nonprincipal ultrafilter
$\mathfrak u$ on $\mathbb N$ and, for every $g\in G_\Q$,
define
\[
                         s(g)=\lim_{n\to\mathfrak u}s_n(g)
                              \quad\hbox{in }E.
\]
Compact Hausdorffness gives this limit. Continuity of
multiplication in $E$ proves
$s(gh)=s(g)s(h)$; continuity of $E\to G_\Q$ gives
$\operatorname{pr}(s(g))=g$. Thus $s$ is an abstract section.
This argument does not assert continuity of a limit of
continuous maps with the pointwise topology.

We prove exactly the local compatibility required below.
Characteristic-zero geometric base change identifies
$\pi_1(U_{\Q_3})$ with $E\times_{G_\Q}G_{\Q_3}$,
so the local point section $s_a$ is viewed as a homomorphism
$G_{\Q_3}\to E$.
For a finite arithmetic quotient $\varphi:E\twoheadrightarrow F$,
let $P=\varphi(\Pi)$. Its finite \emph{\'etale} Galois
cover $W\to U$ has constant quotient
$W/P=U\times\Spec K_0$, where $K_0/\Q$ is the finite
constant extension with group $F/P$. This is the ordinary
arithmetic fundamental-group covering correspondence
\citep[Expos\'e~IX, \S6]{E_SGA1}.
For $b$ sufficiently close to $a$ in $U(\Q_3)$, Hensel's
lemma gives a Galois-equivariant, $F$-equivariant bijection
\[
                         W_b(\overline\Q_3)
                              \xrightarrow{\sim}
                         W_a(\overline\Q_3)
\]
inducing the identity on the indicated constant quotient.
Here is the uniformity in that statement. Choose a finite
Galois extension $L/\Q_3$ splitting the finite fibre $W_a$.
The \emph{\'etale} inverse-function form of Hensel's lemma
gives, through every point of that fibre, a unique local
$L$-analytic sheet over a sufficiently small ball about $a$.
There are finitely many sheets, deck transformations, and
elements of $\Gal(L/\Q_3)$, so one common ball works.
Uniqueness of these sheets makes all their actions compatible.
The constant coordinates remain their original algebraic
values. This proves the asserted equivariance and the identity
on $W/P$.

In torsor coordinates, such a bijection is multiplication by
an element of $P$, since it is the identity modulo $P$.
Consequently for all sufficiently large $n$ the restrictions
$\varphi s_n|_{G_{\Q_3}}$ and
$\varphi s_a|_{G_{\Q_3}}$ are conjugate by $P$.
This is geometric conjugacy, not merely conjugacy by an
arbitrary arithmetic element of $F$. The $P$-conjugacy orbit
of the latter homomorphism is a finite closed subset of
$F^{G_{\Q_3}}$. Its ultralimit therefore still belongs to
that orbit. For each $\varphi$ there is a geometric
conjugator at this finite level.

The set of elements $\gamma\in\Pi$ satisfying
\[
 \varphi\big(\gamma s(g)\gamma^{-1}\big)=\varphi(s_a(g))
                    \qquad(g\in G_{\Q_3})
\]
is nonempty and closed in $\Pi$. Any finitely many such
conditions have a common refinement given by a finite
quotient of $E$. Compactness of $\Pi$ supplies one
$\gamma$ satisfying all of them. Finite quotients separate
elements of $E$, so
\[
               \gamma s|_{G_{\Q_3}}\gamma^{-1}=s_a
\]
as actual homomorphisms. In particular the local restriction
of the abstract section is continuous.

Its global discontinuity is independently visible in one
finite quotient. The cover $z^2=t$ gives a continuous character
$\tau:E\to\F_2$ such that
$\tau s_n=\kappa_2(x_n)$. Independence of the square classes
gives a continuous surjection
\[
 \kappa:G_\Q\twoheadrightarrow\prod_{n\geq1}\F_2,
                 \qquad g\longmapsto(\kappa_2(x_n)(g))_n.
\]
Indeed it is surjective on every finite product by Kummer
theory, and compactness gives surjectivity onto the product.
The ultrafilter functional
\[
 \epsilon:\prod_n\F_2\longrightarrow\F_2,
              \qquad (v_n)\longmapsto\lim_{n\to\mathfrak u}v_n
\]
is additive. It kills the dense direct sum but sends the
constant-one vector to one, so it is discontinuous. A
continuous surjection of compact Hausdorff groups is a
quotient map; hence $\epsilon\kappa$ is discontinuous.
By construction $\tau s=\epsilon\kappa$, proving the
claimed discontinuity without any section-conjecture input.

Realize $E$ as the arithmetic automorphism group of the
universal pro-cover of $U$. Each $s(g)$ is a semilinear
automorphism, and the homomorphism law gives the required
abstract action. It normalizes $\Pi$, which is normal in
$E$. At every finite covering field the automorphism extends
uniquely to its proper normalization; compatibility gives
the action on the completed pro-cover. The stabilizer of the
function $\sqrt t$ is the nonopen subgroup
$\ker(\epsilon\kappa)$, so the induced action itself is
discontinuous. The explicit sequence and all finite
approximations are effective; choosing a nonprincipal
ultrafilter is a choice-principle step and does not give an
algorithm for the abstract section.

In the crossed-cocycle form requested in R22, fix a continuous
reference section $\rho_0$ at $1/2$ and put
\[
 a_g=s(g)\rho_0(g)^{-1}
     =\lim_{n\to\mathfrak u}s_n(g)\rho_0(g)^{-1}\in\Pi.
\]
Then
\[
 a_{gh}=a_g\,\rho_0(g)a_h\rho_0(g)^{-1},
                 \qquad a_g\rho_0(g)=s(g).
\]
Thus this explicitly specifies an abstract crossed cocycle
for the reference action, and its twisted action is exactly
the action constructed above. D.14 computes its fixed set
as empty.
\end{proof}

\begin{claim}{D.13}{Proved}
Let $L/\Q_p$ be finite, let
$\chi:G_L\to\widehat\Z^\times$ be the full cyclotomic
character, and put $H=\ker\chi$. If $b\in L^*$ satisfies
\[
            \kappa_p(b)|_H=0
                    \quad\hbox{in }H^1(H,\Z_p(1)),
\]
then $b$ is a root of unity. In particular, if the restrictions
of the point sections of $\G_m/L$ at $b_1,b_2\in L^*$
are geometrically conjugate on $H$, then $b_1/b_2$ is a
root of unity.
\end{claim}
\begin{proof}
Choose compatible $p$-power roots and let
$c:G_L\to\Z_p(1)$ be the resulting continuous Kummer
cocycle. Since $H$ acts trivially on the coefficient module,
every $H$-coboundary is zero. The hypothesis therefore says
$c|_H=0$ literally. The cocycle descends to a continuous
cocycle on the compact abelian quotient
$A=\chi(G_L)$: its values are constant on $H$-cosets,
and the compact quotient topology gives continuity. This
proves the needed inflation assertion directly for the
compact coefficient module.

There is $a_0\in A$ with $\chi_p(a_0)\ne1$, since a finite
extension of $\Q_p$ does not contain all $p$-power roots
of unity. The cocycle identity, evaluated on $a_0a=aa_0$,
gives
\[
 (\chi_p(a_0)-1)c(a)=(\chi_p(a)-1)c(a_0)
                                         \qquad(a\in A).
\]
The right side is a coboundary. Thus
$(\chi_p(a_0)-1)[c]=0$ in $H^1(A,\Z_p(1))$.
Writing the nonzero scalar as $p^m$ times a $\Z_p$-unit
shows that $p^m\kappa_p(b)=0$ after inflation to $G_L$.
This is the elementary cocycle calculation often called
Sah's lemma; no vanishing assertion for a nonabelian
quotient is being used.

It follows at every finite Kummer level that
$b^{p^m}\in (L^*)^{p^n}$ for every $n$. We record the local
intersection precisely. Write
\[
 L^*=\varpi^{\Z}\times\mu_{q-1}(L)\times U^1,
               \qquad U^1=1+\mathfrak m_L,
\]
where $q$ is the residue-field cardinality. Hensel's lemma
gives the displayed prime-to-$p$ roots of unity. The group
$U^1$ is pro-$p$, since its quotients modulo
$1+\mathfrak m_L^r$ have $p$-power order. An element
divisible by every $p^n$ in $L^*$ has valuation zero;
after removing its unique $\mu_{q-1}$ part, its image in
each finite $p$-group quotient of $U^1$ is a $p^n$-th
power for every $n$, and therefore identity. The filtration
separates points. Hence
\[
                    \bigcap_{n\geq1}(L^*)^{p^n}
                             =\mu_{q-1}(L).
\]
Thus $b^{p^m}$ is a root of unity, and so is $b$.

The final assertion is the Kummer description of
$\pi_1(\G_m/L)=\widehat\Z(1)\rtimes G_L$ with the
splitting at $1$. Point-section classes are Kummer classes
and their difference is the class of $b_1/b_2$
\citep[\S2.2 and \S8.1--8.3, especially Proposition~26]{F_Stix2008}.
Restrict to $H$ and apply the result just proved.
\end{proof}

\begin{claim}{D.14}{Refuted}
The universal assertion D5, that every abstract semilinear
$G_\Q$-action on the universal pro-cover has exactly one
$G_\Q^\circ$-fixed completed geometric closed pro-point,
is false. The action of D.12 has \emph{no} such fixed
pro-point, even though it normalizes the original geometric
deck group.
\end{claim}
\begin{proof}
Here $G_\Q^\circ=\ker\chi=G_{\Q(\mu_\infty)}$.
The completed pro-points are exactly the compatible systems
of geometric closed points on the proper finite levels,
as in D.7; the generic scheme point is not part of this
set. Suppose that a completed pro-point $\xi$ were fixed
by $s(G_\Q^\circ)$, and let
$x\in\mathbb P^1(\overline\Q)$ be its image at the
base compactification.

Use the fixed embedding $\overline\Q\hookrightarrow
\overline\Q_3$ underlying the decomposition group in D.12,
and choose a finite $L/\Q_3$ containing the coordinates
of $x$. Put
\[
                 H=G_L\cap G_\Q^\circ
                      =\ker(\chi:G_L\to\widehat\Z^\times).
\]
Conjugate $s$ by the geometric element supplied by D.12,
and conjugate $\xi$ with it. Its base point remains $x$,
and we now have $s|_{G_{\Q_3}}=s_a$.
Consequently $s_a(H)$ fixes a completed pro-point over $x$.

Choose the following unit of the open tripod:
\[
\begin{array}{c|c|c}
 x&f(t)&f(x)\\\hline
 x\in U(\overline\Q)&t&x\\
 0&1-t&1\\
 1&t&1\\
 \infty&t/(t-1)&1
\end{array}
\]
In every case $f$ is regular and nonzero at $x$.
The morphism $f:U\to\G_m$ therefore sends the decomposition
group of the fixed pro-point over $x$ into the point
decomposition group at $f(x)$. At a cusp its geometric
inertia is killed, since the order of $f$ there is zero.
Equivalently, the covers $w^{3^n}=f(t)$ extend
\emph{\'etale} across $x$ and their fibre at $x$ consists
of the $3^n$-th roots of $f(x)$. Evaluating the section
in these covers gives
\[
       \kappa_3\bigl(f(a)/f(x)\bigr)|_H=0
                         \quad\hbox{in }H^1(H,\Z_3(1)).
\]
This is the usual functorial evaluation of units on
sections \citep[Proposition~26 and \S8.3]{F_Stix2008}.
Changing the selected roots changes only a coboundary;
on $H$ that coboundary is zero because the cyclotomic
action is trivial. The argument applies to interior
points and to each cusp, with no supposition that the
original abstract global section is continuous.

Both $f(a)$ and $f(x)$ belong to $L^*$, and D.13 gives
$f(a)/f(x)\in\mu(L)$. The number $f(x)$ is algebraic
over $\Q$, as is every root of unity. Hence $f(a)$
would be algebraic over $\Q$. Each chosen $f$ is an
invertible fractional-linear function over $\Q$, so
this would make $a$ algebraic, contradicting D.12.
Thus the fixed set is empty.
It is the existence part of the asserted unique fixed point
that fails in this example; no assertion of two distinct
fixed points is made.

D.11 and this counterexample address the two demands
of D5 separately: D.11 constructs an action normalizing
no finite-quotient deck group and having exactly one fixed
cusp, while the present action disproves the universal
fixed-point claim. The latter uses discontinuous abstract
actions and proves no counterexample to the continuous
section conjecture in D1--D4.
\end{proof}
